\subsection{The component-supported OS49 framing and every fan}
\begin{statement}{P3K.34}{Unresolved; scalar-fan audit correction}
\par\noindent\textbf{Audit correction, 17 November 2026.} The component-supported source, line and individual reference characteristic calculation are retained. Section 6 imports the unverified full-plumbing product and P3K.29/P3K.31 recurrence. Its universal fan expression and range are withdrawn. The separate OS49 outer markings require their own audit. The formulas and earlier argument below are retained as the audited prior claim, not as current proved conclusions. See P3K.61--62 and the correction ledger.
\medskip


Retain the actual component-supported OS49 families of P3F.31--32,
with \(\lambda=r^{3k}\), \(k>0\), and cap vector
\(n(1,-1,0,0)^t+d\delta\), where \(n=\pm1\) and \(d\in\mathbb Z\).
For one positive physical \(r\)-loop, in the initial positive-diagonal
Cox product the individual integers and relative coordinates are
\[
 k_{\rm Jin}=-6k,\qquad \ell_{\rm Jin}=-24k,\qquad
 (m_{\rm rel},n_{\rm rel},s)=(6k,0,0).
 \tag{OC.1}
\]
The pair is the same for both norm signs. The line used has its actual
\(I_2\) multidegree \((-1,+1)\) and relative square \(-2k\).
For every actual allowed index-one fan \(T(v,w)\), write the physical
delta cocharacter as \(\delta=a v+b w\), and set \(S=a^2-ab+b^2\).
Its oriented return class is
\[
 [\Sigma_{\Phi,T}]=\frac{3k}{4}(S-1)(S-3)\pmod{28}.
 \tag{OC.2}
\]
At fixed \(k\), its exact range is
\(k\{0,4,6,8,14,18,20,22\}\pmod{28}\).
These conclusions retain the component line, both signed norm maps,
the simultaneous closed Cox product, and the actual common-exterior
frame comparison. They use the stated physical parameter loop.

\end{statement}
\begin{proof}[Preserved earlier argument; scope corrected above]
\par

\medskip\noindent\textbf{1. Literal origins and the simultaneous bounded Cox product}

\par\nobreak

We use P3F.32's actual periods, not a relabelled scalar parameter:
\[
 u=(2s,s)^t+U(q),\quad
 \gamma=(-3s,-s+mv)^t+C(q),\quad D=3U+2C,
\]
\[
 H=-u-\gamma,\qquad K=3u+2\gamma=D+(s+2mv)\delta,
 \qquad \lambda=r^m.
 \tag{OC.3}
\]
Here \(s=\log q/(2 \pi i)\) and \(v=\log r/(2 \pi i)\). In the actual physical
cocharacter lattice the tropical vectors of \(H,K\) are \((1,0),(0,1)\).
The period lattice generated by \((2,1),(-3,-1)\) has determinant one.
On the auxiliary \(2m\)-root the two actual origins are
\[
 \sigma=nvD+nmv(v-1)\delta,\qquad
 o=-\frac{nsD}{2m}-\frac{ns^2}{4m}\delta-nsv\delta.
\]
Their differences are exactly
\[
 (\sigma-o)(v+\theta)- (\sigma-o)(v)
       =n\theta K+nm\theta(\theta-1)\delta,
\]
\[
 (\sigma-o)(s+2m,v)- (\sigma-o)(s,v)=nK+nm\delta.
 \tag{OC.4}
\]
The ordinary closing is multiplication by \(t^n=q^{-n}\), and
\(\sigma(v+1)-\sigma(v)=nK-ns \delta\); hence the actual ordinary endpoint
minus the cusp-origin endpoint is the integral period \(-nK\). This
retains the supplied source-root descent and all endpoint scalar factors.

Put \(U_K=\exp(2 \pi i D_1)\) and \(W_K=\exp(2 \pi i D_2)\). At an initial
parameter \(\lambda_0=r_0^m\), the bounded angular map determined by this
origin is literally
\[
 C_\theta:\quad
 (Z,V)\longmapsto
 \left(U_K^{n\theta}Z,
       (W_K\lambda_0^2)^{n\theta}
       e^{2\pi i nm\theta(\theta-1)}V\right).
 \tag{OC.5}
\]
All powers use the actual logarithms of the nonzero units. At the
endpoint this is the bounded \(nK\) period after reindexing the cell by its
integral vertical shift. The scalar quadratic factor is the actual
physical delta action. We keep it as a circle gauge; in the orbit normal
three-plane it acts identically.

For clarity, this product closes on all three curve charts at once.
In the four-coordinate Cox model of K17, after the period-cell reindexing,
the return has multipliers
\[
 (z_0,z_4,z_1,z_2)\longmapsto
 ((ab\lambda_0^{2n})^{-1}z_0,z_4,az_1,
                     b\lambda_0^{2n}z_2),
 \tag{OC.6}
\]
where \(a,b\) are the retained signed bounded period units. The product
of the four displayed coordinate multipliers is one, so the same map fixes \(q=z_0z_4z_1z_2\) in every chart. The
period lattice acts freely and transitively on cells. Thus no residual
cell permutation or finite stabilizer is hidden in (OC.6). The same
interval of multipliers as in (OC.5), including any odd closed Cox loop,
is used on the two triple charts and all three curve tubes. The finite
cap's additional character \(\exp(2 \pi i d/m)\) is fixed in \(r\); it adds no
parameter difference and remains part of the line gluing in Section 5.

\par

\medskip\noindent\textbf{2. The mixed boundary coordinates and the exact Thom sign}

\par\nobreak

We distinguish a cycle from its point angle. Let \((x,y)\) be the point
angles in the ordered cycles \((u,\gamma)\), and \((h,\kappa)\) those in
\((H,K)\). Equation (OC.3) gives the actual torus isomorphism
\[
 \binom h\kappa=
 \begin{pmatrix}2&-3\\1&-1\end{pmatrix}\binom xy,
 \qquad
 \binom xy=
 \begin{pmatrix}-1&3\\-1&2\end{pmatrix}\binom h\kappa.
 \tag{OC.7}
\]
The other source cycle is \(w\); it is unchanged. These are pointwise
torus maps, with unimodular matrices, in the actual incoming collar.

The marked finite-cap map of the signed page calculation applies to
this cap without changing its elliptic or line action. Its norm is
\(m y-n \alpha\), where \(\alpha\) is the finite base meridian. The explicit
solid-torus map (SCP.1) therefore gives the positive weighted Hopf
angle \(\beta=n y\). Its sheet-zero graph is the actual holomorphic root
slice for either sign, so the page has its complex orientation and
boundary orientation \(\operatorname{or}(E_u)\,d\beta\).

Let the incoming compression graph supplied by F31 be
\(\kappa=g(h,w)\), where \(g\) is a globally real function. The actual positive
parameter loop translates this graph by \(n \theta\) in the K direction.
Its inverse angular coordinate is
\[
 \theta=n(\kappa-g(h,w)).
 \tag{OC.8}
\]
Combining (OC.7) with \(y=n \beta\), the boundary map is the linear torus
isomorphism
\[
 \binom h\theta=
 \begin{pmatrix}2&-3n\\ n&-1\end{pmatrix}\binom x\beta
                  -\binom0{n g(h,w)}.
 \tag{OC.9}
\]
Its determinant is \(-2+3n^2=1\), for both signs; the graph shear also
has determinant one. The \(u\) and \(H\) source orientations agree with
their complex source orientations: their Tate imaginary periods have
leading terms \(2 \operatorname{Im}(s)\) and \(\operatorname{Im}(s)\), respectively, both positive, and
the retained bounded units cannot reverse this sign on the chosen
small collar. Keeping \(w\) first, (OC.9) proves the literal orientation
identity
\[
 \partial V_m\longrightarrow E_H\times S^1_\theta
                  \quad\hbox{has degree }+1
 \tag{OC.10}
\]
for \(n=+1\) and \(n=-1\) separately.

Orient the actual real compression by \(\partial W=F-E_H\). Then the
oriented product \(W\times S^1_\theta\) attaches to the complex graph page and
gives \(\partial X=F\times S^1_\theta\). In the fixed normal chart the delta
weights are \((-1,+1)\). The direct Hopf computation below gives
\(\operatorname{or}(Y)=\operatorname{or}(F)\operatorname{or}(\nu_{\rm Hopf})\). At the outgoing trace the tangent orientation
is \(\operatorname{or}(F)\,dr\,d\theta\); moving the three normal vectors past these two
parameter vectors has positive sign. Hence the normal meridian has
the positive physical Hopf orientation for both norms, and the actual
Euler line satisfies \(\pi_* x=+1\). The two bordism boundary reversals
and Jin's inverse mapping-torus convention are exactly the literal
ones of K7. The complex page and its characteristic numbers therefore
enter with the same sign in both cases.

This calculation is needed even though F32's index is zero. We do not
use multiplication of that zero by \(n\) to determine either integer.

\par

\medskip\noindent\textbf{3. The decreasing K height and the initial Cox normal frame}

\par\nobreak

At the fixed vertical edge the actual positive-fan charts are
\[
 (A,B,Z)=(q/V,V/Z,Z),\qquad
 (A',B',Z')=(qZ/V,V,1/Z),
\]
\[
 A'=ZA,\quad B'=ZB,\quad Z'=Z^{-1},\qquad q=ZAB.
 \tag{OC.11}
\]
The physical delta action has weights \((-1,+1,0)\). Put
\(y=AB\) and \(b=(|B|^2-|A|^2)/2\). Its exact quotient on the two moving
coordinates is \((y,b)\). On the regular collar increasing \(\log|V|\)
increases \(b\), whereas positive \(K\) multiplies \(V\) by
\(q \lambda_0^2 W_K\), whose modulus is less than one. Thus the proper
K-height \(\chi_K\) has \(\partial\chi_K/\partial b<0\). The actual relative orbit product
of F31 is chosen with this collar marking. Its two local orbit
coordinates have positive transition derivative, as in F31/F9, so the
decreasing choice extends through both nodal charts. We retain this
negative derivative.

The initial global quaternion normal coordinate is
\[
 q_H=(1+jZ)(A+j\overline B)
     =(Z^{-1}+j)(A'+j\overline {B'}).
 \tag{OC.12}
\]
The equality follows by multiplication using \(zj=j\bar z\).
At \(Z=0\), the map from the complex normal plane \((A,B)\) to
\(A+j\overline B\) preserves real four-dimensional orientation.
Physical delta is right multiplication by \(\exp(-i \theta)\). Its
positive Hopf normal coordinate is \(\eta=q_Hi\overline{q_H}\).
At \(Z=0\) this is explicitly
\[
 \eta=-2b\,i+2\operatorname{Im}(y)j-2\operatorname{Re}(y)k.
 \tag{OC.13}
\]
The determinant from \((\operatorname{Re} y,\operatorname{Im} y,b)\) to \((\eta_i,\eta_j,\eta_k)\) is \(-8\).
In the circle-first convention, a direct evaluation at \(A=1,B=0\)
has determinant \(-1\) for the ordered coordinates
\((\delta,Z_{\rm re},Z_{\rm im},y_{\rm re},y_{\rm im},b)\) relative to the complex orientation.
The two signs give exactly \(\operatorname{or}(Y)=\operatorname{or}(F)\operatorname{or}(\eta)\). These equations
also identify the orientation at every point by the actual global
quaternion coordinate (OC.12).

Use the oriented raw normal coordinates \((\operatorname{Re} y,\operatorname{Im} y,-b)\) at the
initial chart. Their change to (OC.13) has positive determinant.
Modulo the fixed-curve tangent and its graph shear, the actual
K-product derivative in these coordinates is
\[
 Q_K^+=\begin{pmatrix}
 1&0&0\\0&1&0\\
 (\chi_K)_{\Re y}&(\chi_K)_{\Im y}&-(\chi_K)_b
 \end{pmatrix},\qquad-(\chi_K)_b>0.
 \tag{OC.14}
\]
The source shears contract linearly and the positive last entry
contracts to one. The splittings of the intrinsic K normal line and
the source normal plane form a contractible space, so these
contractions agree on the two charts. They are independent of the
parameter angle. This supplies the initial Cox-to-product normal
isomorphism with its actual reflection retained at both endpoints.

For the signed rotating coordinate \(n \kappa\), the raw derivative is
\(T_n Q_K^+\), \(T_n=\operatorname{diag}(1,1,n)\). We apply the same constant coordinate
change \(T_n\) to both initial normal frames. Its remaining relative
comparison is \(T_n^{-1}(T_n Q_K^+)=Q_K^+\). Equivalently its conjugation
on a full clutch is conjugation by \(n T_n\), an element of \(SO(3)\).
Thus neither the decreasing K-height nor \(n=-1\) supplies an additional
free three-dimensional framing term.

We also retain the normal derivative of the actual bounded map (OC.5).
A unitary torus multiplier \((Z,V)\longrightarrow (aZ,bV)\) acts in (OC.11) by
\((Z,A,B)\longrightarrow (aZ,pA,q_1B)\) with \(p=b^{-1}\), \(q_1=b/a\), and \(a p q_1=1\).
Choose continuous roots along its given parameter interval,
\(L^2=p q_1\), \(R^2=p/q_1\), \(LR=p\). Direct substitution gives
\[
 q_H\longmapsto L q_H R,\qquad
 \eta\longmapsto L\eta\overline L.
 \tag{OC.15}
\]
The latter rotation is constant along the fixed sphere. Positive real
factors are the same normal isomorphism in the two products; they
cancel in the relative comparison. Higher terms of the retained base
units have zero normal linear part because \(q=ZAB\) is quadratic in
the normal coordinates. Consequently the remaining closed rotation
loop, including an odd one, extends over the fixed three-disc. The
actual simultaneous Cox product has the same free normal-framing
class as the rotating product. No endpoint loop has been replaced by
a chosen trivial loop.

\par

\medskip\noindent\textbf{4. The degree-2m graph phase and the full normal clutch}

\par\nobreak

The actual one-cap weighted page has regular-binding polynomial
\(z_2-t_* z_1^m\); its source base is \(t=z_2/z_1^m\).
Thus in the actual holomorphic binding chart \(q_b=t-t_*\), after
retaining its fixed phase, the page equation is
\(\alpha=\phi-m \beta\). If the cusp coordinate is \(q=1/t\), this same
chart is \(q_b=-(q-q_*)/(q q_*)\), a specified holomorphic unit times
\(q-q_*\); it has nonzero derivative. We use this exact chart throughout
the collar calculation. F31's norm-shell interpolation keeps the base
coordinate fixed and hence preserves this equation.

On the rotating incoming graph, (OC.7)--(OC.8) give
\[
 \beta=n(-h+2\kappa)=2\theta+g_n(h,w),\qquad
 g_n=n(2g-h),
\]
\[
 \alpha=\phi-2m\theta-m g_n(h,w).
 \tag{OC.16}
\]
The function \(g_n\) is circle-valued; its \(-n h\) term is retained.
Since \(m\) is integral, \(R_z(-m g_n)\) is a global map. The factor two
is forced by the actual mixed point-angle map, and is not a further
parameter cover.

Use the relative angle
\(\eta_{\rm rel}=n(\kappa-g(h,w))-\theta\), which vanishes on the graph.
After quotienting by the source tangent, its two parameter Cartesian
columns have the same relative coefficients \((-\sin \theta,\cos \theta)\)
as in K30. The derivatives of \(g\) are precisely the removed triangular
source shears. Give the graph its oriented page-normal vector and the
Cartesian parameter frame. The common outgoing page and incoming
compression collar has tangent
\((\cos \alpha,\sin \alpha,\rho,0)\), \(\rho>0\), in the binding-plane,
parameter-radius, relative-angle quotient. The exact quotient map is
\[
 (v_q,v_r,v_\eta)\longmapsto
 \left(v_q-\frac{v_r}{\rho}(\cos\alpha,\sin\alpha),v_\eta\right).
 \tag{OC.17}
\]
Projecting the two Cartesian parameter vectors and retaining the same
positive reference rescalings gives the actual matrix
\[
 G(\alpha,\theta)=
 \begin{pmatrix}
 -\sin\alpha&\cos\theta\cos\alpha&\sin\theta\cos\alpha\\
 \cos\alpha&\cos\theta\sin\alpha&\sin\theta\sin\alpha\\
 0&\sin\theta&-\cos\theta
 \end{pmatrix}\in SO(3).
 \tag{OC.18}
\]
The constant signed/reflected coordinates in Section 3 are applied
at both ends. Equation (OC.10), followed by the local calculation
(OC.13), fixes the positive Hopf normal orientation of (OC.18).

Choose a normal frame on the compression at parameter zero and
transport it by the actual K flow with speed \(n\). Its difference at
the fixed sphere from (OC.12) is a map \(S^2\longrightarrow SO(3)\) and extends over
the fixed three-disc. The construction is possible because the
compression is a solid torus minus a ball, its tangent bundle is
trivial, and the actual \(Y=S^5\) is stably parallelizable; the oriented
rank-three normal bundle has zero \(w_2\) and hence admits a frame on
this three-dimensional body. In the transported product the frame is
\(R(p)\), independent of \(\theta\).

Since \(G(\alpha,\theta)=R_z(\alpha)G(0,\theta)\), the full actual clutch is
homotopic, by the explicit shear and positive collar homotopies, to
\[
 \begin{split}
 R(p)^{-1}G(\phi-2m\theta-mg_n(p),\theta)
   ={}&[R(p)^{-1}R_z(-mg_n(p))]\\
     &\cdot[R_z(\phi-2m\theta)G(0,\theta)].
 \end{split}
 \tag{OC.19}
\]
The first factor comes from the actual source \(T^2\), the second from
the physical parameter circle. They are globally defined even when
the graph has winding or the frame has a nontrivial one-dimensional
class. For the product of two such maps, the Maurer--Cartan cubic
trace form is the sum of the factor cubic forms plus an exact form;
this follows by expanding \((ab)^{-1}d(ab)\) and using the two
Maurer--Cartan equations. Each factor cubic form vanishes by
dimension. The interpolating connection on the clutch collar
therefore gives
\[
             \langle p_1(\overline\nu),[\overline X]\rangle=0.
 \tag{OC.20}
\]
This evaluates the actual normal bundle; it does not assert that all
its lower-degree data vanish.

\par

\medskip\noindent\textbf{5. The component line and the individual integers}

\par\nobreak

F31--32 construct the page by the literal line subfan with rays
\((j,\lceil j/2\rceil,1)\), whose two source component degrees are \((-1,+1)\).
Its m separated norm shells embed because the displayed height
defects prevent reentry after the first zero. The positive radial
annular map commutes with the original contracting line lift and
retains its circle line. The finite root and its fixed scalar
character are also the actual ones of F29/F32. Consequently their
compact page and relative line have exactly
\[
 \chi(\mathcal O_{S_{\rm page}})=k,\qquad
 \sigma(S_{\rm page})=-8k,\qquad
 c_1(M)_{\rm rel}^2=-2k,
 \qquad p_{\rm rel}=-24k.
 \tag{OC.21}
\]
The \(m\) regular half-disc bands of F32 prove equality of the graph
meridian line with this actual line: their relative \(H^2\) vanishes, so
restriction is injective. The fixed \(d\) character is a base line on
the page; after \(m\) sheets its character is one, and an integral fibre
class change has square and pairing with \(M\) zero. All nonzero \(I_2\)
component degrees remain in the computation (OC.21).

The two cap markings are the literal unimodular map (OC.7); their
closed difference is \(S^1_u\times S^3\). The line and signature of that
difference are zero by the actual cap product maps in F31. Thus
the signed trace calculated in Sections 2--4 has
\[
 b^2=4c_1(M)_{\rm rel}^2=-8k,\qquad
 p_1(T\overline X)=-24k,
\]
with \(\pi_*x=+1\) for both signs. The sphere-bundle pushforward equations
of K7, using the matched initial Cox product, give
\[
 k_{\rm Jin}=\frac{p_1(\overline\nu)+3b^2}{4}=-6k,
 \qquad
 \ell_{\rm Jin}=p_1(T\overline X)+p_1(\overline\nu)=-24k.
 \tag{OC.22}
\]
The exact relative coordinate morphism of K23 is
\(m_{\rm rel}=-\ell/4\), \(n_{\rm rel}=-(4k_{\rm Jin}-\ell)/4\),
\(s=(4k_{\rm Jin}-\ell)/24\pmod{28}\). Hence (OC.1) follows. This calculation
uses the component line itself and does not infer the pair from the
already zero index.

\par

\medskip\noindent\textbf{6. The three-curve reference and all actual fans}

\par\nobreak

The simultaneous product (OC.5)--(OC.6) gives a delta-equivariant
representative of the physical return equal to the identity near
all three cusp curves. The period lattice is index one, so the
common-exterior dual cells, primitive circle-trace kernel, and
complex-meridian maps of K29 apply to these actual charts unchanged.
In particular, the variation is
\[
       V=\operatorname{diag}(-6k,0,0).
\]
We recall why the two nonfixed stable derivatives vanish here. On
the free circle complement an invariant connection expresses the
stable derivative as the pullback of the quotient derivative, with
a contractible off-diagonal connection block. The two corresponding
relative generators are exactly the actual delta traces of integral
two-cycles. A pulled-back third obstruction evaluates to zero on
them. Their variations are zero, so the exact frame-change formula
\(c_{\sigma g}(F)=g(Fp)^{-1}c_\sigma(F)g(p)\) changes those values by zero.
Thus their K23 relative derivative integers are zero in the actual
single-curve product frames. On the fixed curve (OC.1) gives
\(n_{\rm rel}=0\); K31's literal first-resolution product-frame change gives
the initial ambient derivative vector
\[
                  \delta_{\rm der}=(-6k,0,0)^t.
 \tag{OC.23}
\]

For an actual periodic flop of edge \(i\), K29/K31 give the common-exterior
maps
\[
 V'=M_iVM_i^t,\qquad
 \delta'_{\rm der}=M_i(\delta_{\rm der}-2Ve_i),\qquad
 s'=s+(\delta_{{\rm der},i}-V_{ii})/2.
 \tag{OC.24}
\]
Here \(M_i\) has column \(i\) equal to \(-e_i+2\sum_{j\ne i}e_j\) and other
columns unchanged. These are actual two blowdowns, their stable
frame transition of degree \(+2\), and the orientation-preserving
anti-quaternion reset, all intertwining the same Cox parameter
product. They do not depend on the whole-fibre assumption on a
line elsewhere. We have supplied the needed closed product and
initial derivative data for the present component-supported line.

Write \(V=-6k ww^t\), initially \(w=(1,0,0)\), and
\(\delta_{\rm der}=a_{\rm der}w\), initially \(a_{\rm der}=-6k\). Let \(S=\sum_iw_i\).
The recurrence gives \(w'=M_iw\), \(S'=S+2w_i\) and the exact central
increment
\[
 s'-s=-3k(S'-S)+\frac{3k}{4}(S'^2-S^2).
\]
It telescopes to (OC.2). P3F.35 connects every allowed one-vertex
unimodular periodic fan to the initial fan by these literal moves,
preserving the actual smooth sphere and all exterior data. In the
triangle with rays \((0,1),(v,1),(v+w,1)\), the physical delta weights
are \((-a,a-b,b)\), so
\(S=(a^2+(a-b)^2+b^2)/2=a^2-ab+b^2\). Every primitive pair \((a,b)\) is
realized by the explicit integral basis map of K31. This proves the
all-fan formula with the fan marking retained. The two diagonal
fans retain the physical vertical edge and have \(S=1\), giving zero,
as independently proved for those fillings in F32/K17.

For completeness, its exact range at fixed \(k\) is
\[
       k\{0,4,6,8,14,18,20,22\}\pmod {28}.
 \tag{OC.25}
\]
For a primitive pair \(S\) is odd, so \((S-1)(S-3)\) is divisible by eight.
The value in (OC.2) is determined by \(S\) modulo 56. Equivalently we
enumerate pairs modulo 56 with \(\gcd(a,b,56)=1\). Each such residue pair
has a primitive integral lift: choose a nonzero lift of \(b\) and use
the Chinese remainder map to choose \(a\) modulo 56 avoiding its finitely
many prime divisors not dividing 56. The remaining prime divisors
already do not divide the specified \(a\) residue. The exact enumeration
and primitive witnesses are in the certificate. For \(k=1\), witnesses
\[\begin{gathered}(a,b)=(0,1),(1,10),(1,4),(3,14),\\(1,6),(1,3),(2,5),(3,7)\end{gathered}\] give the
listed values in that order. This enumeration assigns no value to
an unconstructed fan: K31's explicit basis matrix constructs every
one of these fans in the actual period lattice.

\par

\medskip\noindent\textbf{Certificate and scope}

\par\nobreak

The exact origins, mixed coordinate matrices for both signs, the Hopf
orientation determinant, reflected normal comparison, quaternion
torus map, degree-2m normal factorization, characteristic pair, and
primitive fan residue enumeration are executed in
\path{code/part3/os49_full_framing.py}, with JSON and captured output.
The proof is for the positive physical loop in \(\lambda=r^{3k}\);
there is no division by the page degree or by a scalar cover. The
fixed \(d\) character is retained, and no additional parameter-dependent
clutch is introduced.

\end{proof}
